% ch1_answers.tex -- answers to the Quick check questions in Module 7, Chapter 1.
% (No nested lists here: every \item is one answer.)
\begin{chapteranswers}{1}{The Electromagnetic Spectrum}

% ---- 1.1 Maxwell and EM waves
\item Maxwell added: \emph{a changing electric field makes a magnetic field}. Faraday's law
  already said that a changing magnetic field makes an electric field. With Maxwell's addition,
  each changing field makes the other. So the two fields can keep each other going in empty
  space, with no charges or currents. This travelling pair is an EM wave. Without the addition,
  a changing electric field in empty space would make no magnetic field, and there would be no
  wave.

\item \eqlines{\mu_0\varepsilon_0 &= (\num{1.26e-6})\times(\num{8.85e-12})
    = \qty{1.1151e-17}{s^2.m^{-2}} \quad\text{(multiply first)}\\
  c &= \frac{1}{\sqrt{\qty{1.1151e-17}{s^2.m^{-2}}}} = \frac{1}{\qty{3.3393e-9}{s.m^{-1}}}\\
    &= \qty{2.9946e8}{m.s^{-1}} \quad\text{(keep the extra digits)}}\\
  Answer: \final{$c = \qty{2.99e8}{m.s^{-1}}$} (3 significant figures, like the data). This is a
  little lower than the \qty{2.998e8}{m.s^{-1}} of \cref{ch1:we:maxwell}. Rounding $\mu_0$ up
  from $1.2566\times10^{-6}$ to $1.26\times10^{-6}$ makes the product $\mu_0\varepsilon_0$
  bigger, so $c$ comes out smaller.

\item (1) An AC supply makes the electrons in the antenna oscillate up and down, so they are
  accelerating. (2) The oscillating charges make a changing electric field, and the changing
  current makes a changing magnetic field. (3) Each changing field makes the other (Faraday's law
  and the Ampère--Maxwell law). So an EM wave, with the same frequency as the AC, spreads out from
  the antenna at $c$.

\item North or south (horizontal). $B$ must be at \ang{90} to $E$ (up) and at \ang{90} to the
  direction of travel (east). The only directions left are north and south. (Extension, not
  needed for the HSC: a right-hand rule links $E$, $B$ and the direction of travel. It gives
  south at this instant.)

\item $f = \qty{702}{kHz} = \qty{7.02e5}{Hz}$. Then
  $\lambda = \dfrac{c}{f} = \dfrac{\qty{3.00e8}{m.s^{-1}}}{\qty{7.02e5}{Hz}} = \qty{427.4}{m}$.
  Answer: \final{$\lambda = \qty{427}{m}$}. AM radio waves are hundreds of metres long.

% ---- 1.2 speed of light
\item The pulse goes there and back. Distance travelled:
  $d = ct = \qty{3.00e8}{m.s^{-1}}\times\qty{1.20e-4}{s} = \qty{3.60e4}{m}$. The aircraft is
  half of this away: $D = \dfrac{d}{2} = \qty{1.80e4}{m}$. Answer:
  \final{$D = \qty{1.80e4}{m}$} (\qty{18.0}{km}).

\item (a) $t = \dfrac{d}{c} = \dfrac{2\times\qty{1500}{m}}{\qty{3.00e8}{m.s^{-1}}}
  = \final{\qty{1.0e-5}{s}}$.\\[4pt]
  (b) Human reaction time is about \qty{0.2}{s}. That is
  $0.2 \div (1.0\times10^{-5}) = \num{20000}$ times longer than the travel time. The two
  people's reaction times hid the tiny delay. The method measured reaction time, not the speed
  of light. His timer also could not measure such a short time.

\item The light must go there and back while the wheel turns from a gap to the next tooth.
  First find the travel time:\\[2pt]
  \eqlines{t &= \frac{2D}{c} = \frac{2\times\qty{8633}{m}}{\qty{3.00e8}{m.s^{-1}}}
    = \qty{5.755e-5}{s}}\\[4pt]
  Rearrange $t = \dfrac{1}{2Nf}$ for $f$. Multiply both sides by $f$: $ft = \dfrac{1}{2N}$.
  Divide both sides by $t$: $f = \dfrac{1}{2Nt}$. Then substitute:\\[4pt]
  \eqlines{f &= \frac{1}{2\times720\times\qty{5.755e-5}{s}} = \qty{12.07}{s^{-1}}}\\[2pt]
  Answer: \final{$f = 12.1$ turns per second}. (Fizeau saw darkness at about 12.6.)

\item The metre is the distance light travels in a vacuum in $\dfrac{1}{299\,792\,458}$ of a
  second, with $c$ fixed at exactly \qty{299792458}{m.s^{-1}}. Any two advantages:
  (1) Any laboratory can reproduce it. Nobody needs to compare with one bar kept near Paris.
  (2) It rests on a constant of nature that does not change. A metal bar changes length with
  temperature, and it can be worn, damaged or lost. (3) It is far more precise, because atomic
  clocks measure time extremely precisely.

% ---- 1.3 spectra
\item Continuous: a hot filament (or lava, or the inside of a star). Emission: a discharge tube
  (or a neon sign, or a flame). Absorption: sunlight, or white light passed through a cool gas.

\item The lines come from electrons jumping between energy levels. In emission, an electron
  drops from a higher level to a lower one and gives out a photon whose energy equals the gap.
  In absorption, an electron takes in a photon with exactly the same gap energy and jumps up.
  The same element has the same gaps. So the photons have the same energies, and therefore the
  same frequencies and wavelengths ($E = hf$ and $c = f\lambda$).

\item Check each element in \cref{ch1:tab:lines}. Helium: 447, 471, 492, 502, 588, 668 \cmark\
  all present. Mercury: 405, 436, 546, 577, 579 \cmark\ all present. Hydrogen: no 656, 486, 434
  or 410 \xmark. Sodium: no 589 \xmark\ (588 belongs to helium).
  Answer: \final{helium and mercury}. Every line is explained, with none left over.

\item Her conclusion is not justified. One line is not enough: another element could have a line
  near \qty{486}{nm}. Hydrogen also has lines at 410, 434 and \qty{656}{nm}. She must check that
  all of them are present. The strong red line at \qty{656}{nm} is the most important. Only then
  can she say that hydrogen is there.

% ---- 1.4 stars
\item $\lambda_\mathrm{max} = \qty{828}{nm} = \qty{8.28e-7}{m}$.
  $T = \dfrac{b}{\lambda_\mathrm{max}} = \dfrac{\qty{2.898e-3}{m.K}}{\qty{8.28e-7}{m}}
  = \qty{3.50e3}{K}$. Answer: \final{$T = \qty{3.50e3}{K}$}. This is a cool star, so it looks
  red.

\item (a) Change to kelvin first: $T = 2527 + 273 = \qty{2800}{K}$.
  $\lambda_\mathrm{max} = \dfrac{b}{T} = \dfrac{\qty{2.898e-3}{m.K}}{\qty{2800}{K}}
  = \qty{1.035e-6}{m}$.\\[4pt]
  Answer: \final{$\lambda_\mathrm{max} = \qty{1.035e-6}{m}$}
  ($=\qty{1035}{nm}$, about \qty{1}{\micro\metre}; the data have 4 significant figures). (b) The peak is in the infrared, beyond \qty{700}{nm}. Most of the energy
  comes out as infrared (heat), and only a small part as visible light. So most of the
  electrical energy is wasted.

\item One edge of the star turns towards us, so its light is blueshifted. The other edge turns
  away, so its light is redshifted. We receive light from the whole star at once, so each line
  is spread over a range of wavelengths: it is broader. The shifts from the two edges are equal
  and opposite. The middle of the disc moves across our line of sight, so its light is not
  shifted. So the line widens evenly on both sides of its usual wavelength, and its centre does
  not move.

\item Star B. The rotation speed of the surface is the same, so rotation cannot explain the
  difference. In a denser
  atmosphere, atoms collide more often and are disturbed by nearby charged particles. Their
  energy levels shift slightly, so each atom absorbs a slightly different wavelength. The lines
  get wider (pressure broadening). So B, with the broad lines, has the denser atmosphere.

\item (a) The star's line is at a \emph{shorter} wavelength than the laboratory line
  ($656.17 < 656.30$). This is a blueshift, so the star is moving \final{towards Earth}. All
  its lines move by the same fraction, so this is a Doppler shift, not a different element.\\[4pt]
  (b) $\Delta\lambda = 656.30 - 656.17 = \qty{0.13}{nm} = \qty{1.3e-10}{m}$. Rearrange
  $\Delta\lambda = \dfrac{\lambda v}{c}$: multiply both sides by $c$, then divide both sides
  by $\lambda$. This gives $v = \dfrac{c\,\Delta\lambda}{\lambda}$. So
  $v = \dfrac{\qty{3.00e8}{m.s^{-1}} \times \qty{1.3e-10}{m}}{\qty{6.5630e-7}{m}}
  = \qty{5.94e4}{m.s^{-1}}$. Answer: \final{$v \approx \qty{5.9e4}{m.s^{-1}}$}, about
  \qty{59}{km.s^{-1}} (2 significant figures, because $\Delta\lambda$ has 2).\\[4pt]
  (c) A Doppler shift shows only the part of the motion along our line of sight (straight
  towards or away from us). Motion across the sky gives no shift. So the star's real speed
  through space can be bigger than \qty{59}{km.s^{-1}}.

\end{chapteranswers}
